\section{E2: tasks where the answer is known in advance}
\label{sec:synthetic}

A null result on CIFAR-10 is ambiguous. It could mean the architecture does not work, or it
could mean CIFAR-10 classes are not concepts that reward a second layer --- in which case the
experiment never had the power to show anything. The two tasks here remove that ambiguity by
being simple enough to analyse on paper.

\lead{Setup.} A $32\times32$ Boolean image containing a \textsc{square} and/or a
\textsc{cross} on $16$ equally spaced slots around a circle, plus $1\%$ salt noise. Layer~1
uses $4\times4$ patches, so \emph{no single patch can ever contain both shapes}.

\lead{Task \code{xor}} --- label $=$ (square present) XOR (cross present). One layer computes
$\mathrm{threshold}\big(\sum_j w_j \bigvee_p \mathrm{conj}_j(\text{patch}_p)\big)$. Its
features are patch-local presence detectors, and ``cross absent everywhere'' is a conjunction
\emph{over} patches that no single clause can form. The class score is therefore a linear
threshold in $(z_A, z_B)$, and XOR is not linearly separable: $w_A > t$ and $w_B > t$ force
$w_A + w_B > t$. Position literals do not rescue it --- the same argument runs through with
$z_{A,\text{top}}, z_{A,\text{bot}}, \dots$ The one-layer ceiling is the best linear
separator, $75\%$. Both stacked variants can express it, \emph{including} the flat stack,
because negating a globally OR-pooled bit does mean ``absent everywhere''.

\lead{Task \code{near}} --- both shapes always present; label $=$ are they close together.
Because the slots lie on a ring, rotational symmetry gives every slot exactly the same number
of near and far partners, so the marginal position of each shape is identical in both classes
and only the \emph{joint} separation carries the label. One layer is out: its additive vote
realises only $f(\mathrm{pos}_A) + g(\mathrm{pos}_B) > t$, and no $4\times4$ patch sees both
shapes. The flat stack is out: it globally OR-pools before the second stage and keeps no
position at all. The spatial \mctm can express it --- a layer-2 clause ``channel $A$ fires
here \emph{and} channel $B$ fires here'' over a $5\times5$ window ($13$ input pixels, enough
for the near separation and not the far one), OR-pooled over positions, is exactly
``somewhere the two are close''.

\begin{ttcallout}[Why the ring matters]
Sampling positions freely and rejecting on distance does \emph{not} give equal marginals: it
pulls near-pairs towards the centre. Fitting the best additive $f(\mathrm{pos}_A) +
g(\mathrm{pos}_B)$ predictor --- precisely the function class one convolutional layer
realises --- reaches $66.2\%$ on that version of the task, without any joint reasoning. On
the ring construction the same predictor reaches $46.6\%$, i.e.\ chance, while the true joint
condition is perfectly predictive. The check is \code{scripts/check\_leakage.py}.
\end{ttcallout}

\lead{The oracle layer 1.} Trained greedily on a task it cannot represent, layer~1 has little
reason to learn clean detectors, so a null result for the stack would again be ambiguous ---
architecture, or training? Each stacked arm is therefore run twice: once with the greedy
layer~1 of the proposal, and once with an \emph{oracle} layer~1 whose two clauses are an exact
square detector and an exact cross detector, written straight into the automata
(\code{build\_oracle\_l1}). Measured recall is $1.000$ for the square and $0.936$ for the
cross, with no false positives, so the oracle's features are correct by construction and
whatever the stack then achieves is a property of the architecture alone.

\begin{table}[h]
\centering\small
\caption{Task \code{xor}: mean best test accuracy over 3 seeds. The one-layer ceiling is
$75\%$ by the argument above.}
\label{tab:synxor}
\pgfplotstabletypeset[
  col sep=comma, string type,
  every head row/.style={before row=\ttoprule, after row=\ttmidrule},
  every last row/.style={after row=\ttbotrule},
  columns={label,acc,sd,seeds},
  columns/label/.style={column name=\thead{arm}, column type={l}},
  columns/acc/.style={column name=\thead{test acc (\%)}, column type={r}},
  columns/sd/.style={column name=\thead{$\pm$}, column type={r}},
  columns/seeds/.style={column name=\thead{seeds}, column type={r}},
]{data/syn_xor.csv}
\end{table}

\begin{table}[h]
\centering\small
\caption{Task \code{near}: mean best test accuracy over 3 seeds. One layer and the flat stack
are both at chance by construction; only the spatial map can carry the joint condition.}
\label{tab:synnear}
\pgfplotstabletypeset[
  col sep=comma, string type,
  every head row/.style={before row=\ttoprule, after row=\ttmidrule},
  every last row/.style={after row=\ttbotrule},
  columns={label,acc,sd,seeds},
  columns/label/.style={column name=\thead{arm}, column type={l}},
  columns/acc/.style={column name=\thead{test acc (\%)}, column type={r}},
  columns/sd/.style={column name=\thead{$\pm$}, column type={r}},
  columns/seeds/.style={column name=\thead{seeds}, column type={r}},
]{data/syn_near.csv}
\end{table}

\subsection{Readout}

\lead{\code{xor} --- depth adds real expressivity.} A single layer is capped at $75\%$ by the
argument above and measures $\synXorSingleTask{} \pm \synsdXorSingleTask{}\%$, with a large
spread: Tsetlin training finds the optimal separator on some seeds and not others. Trained
greedily neither stack escapes the ceiling ($\synXorFlatStackTask{}\%$,
$\synXorMctmTask{}\%$). Given an oracle layer~1 both clear it decisively --- the flat stack
reaches $\synXorFlatStackOracle{}\%$ with two channels and
$\synXorFlatStackOracleSixFour{} \pm \synsdXorFlatStackOracleSixFour{}\%$ with sixty-four,
and the \mctm $\synXorMctmOracle{}\%$ and $\synXorMctmOracleSixFour{}\%$. Because the
sixty-four-channel oracle gives layer~2 exactly the channel count --- and therefore the
literal-space size --- that it faces with a greedily trained layer~1, the gap cannot be read
as ``fewer channels are easier''. A stack of convolutional Tsetlin layers really does express
what one layer provably cannot; the greedy first layer is what stands between the
architecture and that capability.

\lead{\code{near} --- keeping the spatial map is the part of the idea that earns its keep.}
The flat stack is at chance in \emph{every} configuration --- $\synNearFlatStackTask{}\%$
greedy, $\synNearFlatStackOracle{}\%$ with the two-channel oracle,
$\synNearFlatStackOracleSixFour{}\%$ with sixty-four. No quality of features rescues it,
because global OR-pooling has already destroyed the position the task depends on. The single
layer is likewise at chance ($\synNearSingleTask{}\%$). The spatial \mctm reaches
$\synNearMctmTask{} \pm \synsdNearMctmTask{}\%$ with the ordinary greedily trained layer~1.
This is the one place in the whole evaluation where the proposal's distinctive choice is
decisively correct.

\begin{ttkey}
The two tasks answer the two questions separately. \code{xor}: stacking convolutional Tsetlin
layers is strictly more expressive than one layer --- but only greedy training stands in the
way of using it. \code{near}: of the two ways to stack, only the one that keeps the spatial
map can represent a relational concept at all; the cheap global-OR stack is at chance no
matter how good its features are.
\end{ttkey}

\lead{A sparsity result in the other direction.} The two-channel oracle \emph{fails} on
\code{near} ($\synNearMctmOracle{}\%$) while the same oracle padded to sixty-four channels
with random distractors succeeds ($\synNearMctmOracleSixFour{} \pm
\synsdNearMctmOracleSixFour{}\%$). Each oracle channel fires at exactly one location per
image, so a layer-2 clause must pin down one specific pair of offsets, and any given pair is
vanishingly rare in training; the extra channels supply the redundancy that makes a
co-occurrence learnable. Sparsity hurts the second layer from both directions --- too sparse
and conjunctions never fire, too precise and there is nothing to generalise over. It is the
same mechanism that \S\ref{sec:density} measures on CIFAR-10, seen from the opposite side.

